%% P01 \dash{} Floating point: what the machine actually computes
%%
%% Stub. The brief below is the contract for this program: what it must argue,
%% and what it must buy the reader. Writing the program means deleting the
%% \programstub{} block and replacing it with the program. If the brief turns
%% out to be wrong once the mathematics has been worked, change it and record
%% the contradiction in CLAUDE.md under *Resolved questions*.

\program{Floating point: what the machine actually computes}
\label{prog:P01}

\programstub{%
Argument: a float is a sign, an exponent and a fixed number of significant
bits, so the gap between representable numbers grows with magnitude, and
equality is not a question you should ask. Payoff: why \code{bf16} is the
training format and \code{fp16} needs loss scaling \dash{} they have the same
width and different exponent budgets, and the reader computes the overflow
threshold rather than being told it; why a sum of gradients is not
associative and two GPUs therefore disagree in the last bits; what machine
epsilon is and why comparing two loss values that differ in the eighth
decimal is measuring the hardware.

\medskip
\textbf{Planned length:} 45 frames.
}
